\documentclass[10pt,conference]{IEEEtran}
\usepackage[utf8]{inputenc}
\usepackage[T1]{fontenc}
\usepackage{times}
\usepackage{graphicx}
\usepackage{amsmath,amssymb}
\usepackage{booktabs}
\usepackage{multirow}
\usepackage{xcolor}
\usepackage{hyperref}
\usepackage{cite}

\hypersetup{
    colorlinks=true,
    linkcolor=blue,
    citecolor=blue,
    urlcolor=blue
}

\title{IRAG: Insurance Retrieval-Augmented Generation System}

\author{
\begin{tabular}{cccc}
HONG Shiyu & ZHANG Haoran & YANG Kaiwei & ZHANG Bowen \\
59535572 & 60014388 & 59553160 & 60028995
\end{tabular}
}

\begin{document}

\maketitle

\begin{abstract}
This paper presents IRAG, a Retrieval-Augmented Generation system for Hong Kong insurance document question answering, built around a three-stage LoRA fine-tuning pipeline for Qwen3-4B.
Stage~1 grounds the model in clause reading comprehension; Stage~2 adds structured database reasoning via SQL chain-of-thought; Stage~3 applies Retrieval-Augmented Fine-Tuning (RAFT), teaching the model to cite evidence when it is present and refuse when it is not.
Offline evaluation against an untuned baseline shows that fine-tuning improves evidence use (4.07 vs.\ 3.53) and answer compliance (70.7\% vs.\ 62.0\%) while maintaining near-parity on hallucination suppression (83.6\% vs.\ 85.6\%) under the deployment system prompt.
The model achieves 100\% structured output compliance across 15 live pipeline responses, and a grounded attribution judge---introduced to detect hallucinations that standard LLM-as-judge misses---narrows the evaluated hallucination gap from 2.0 to 4.0 percentage points after claim-level source tracing.
End-to-end evaluation on real Hong Kong insurance questions yields a 64.3\% pass rate on answerable queries, with most failures attributable to retrieval-layer gaps rather than model behaviour.
\end{abstract}

\section{Introduction}
% ============================================================

Hong Kong's insurance market is one of the densest in Asia, with over 160
licensed insurers publishing product disclosure documents that mix Chinese and
English terminology, actuarial tables, and legally operative clause text.
Verifying a specific benefit, exclusion, or renewal condition requires locating
the right clause across dozens of PDFs---slow and inaccessible to most
policyholders.
Keyword search breaks down on paraphrased queries; general-purpose LLMs fail
differently, hallucinating policy details, ignoring questions the evidence
cannot support, and producing free-form prose with no audit trail.

The IRAG project tackles this through two successive development phases.
The first phase, carried out in a prior semester, established an initial
Milvus-based RAG prototype: a retrieval pipeline over Chinese insurance text
and a simple table-embedding scheme using a Google text embedding model
designed for English tables.
This prototype identified three concrete limitations.
The knowledge base was narrow in scale and insurer coverage.
The English-oriented table embedder performed poorly on Chinese insurance
tables and represented each table as a single coarse vector, making it
difficult to pinpoint the specific row or clause needed to answer a question.
Pilot fine-tuning experiments---a Chain-of-Thought run and a conventional
insurance-style answer run---showed promise but were not yet integrated with
the retrieval system.

Building on this foundation, the current semester addresses both limitations
in parallel.
On the retrieval side, the team expanded the knowledge base to 247 Hong Kong
insurance product PDFs across nine insurers, replaced the table embedder with
a fine-grained three-modal chunking scheme (full table, columns, rows), added
cross-encoder reranking, and deployed a FastAPI service with a Vue~3 chat
interface.
On the generation side, this paper describes the fine-tuning pipeline designed
to make the model safe to deploy in this RAG context.
We train Qwen3-4B~\cite{qwen3} through three successive LoRA stages: clause
reading comprehension, structured database reasoning via SQL chain-of-thought,
and Retrieval-Augmented Fine-Tuning (RAFT)~\cite{zhang2024raft} that teaches
the model to reason over retrieved chunks and refuse when those chunks do not
contain the answer.
We evaluate the result at three levels---offline behavioral tests against an
untuned baseline, a grounded attribution judge designed to catch hallucinations
that standard LLM-as-judge misses, and end-to-end smoke tests on the live
pipeline---to get a clear picture of what the fine-tuning fixes and where gaps
remain.

The rest of the paper is structured as follows.
Section~\ref{sec:related} reviews related work.
Section~\ref{sec:method} describes the three-stage training pipeline and
the RAG infrastructure.
Section~\ref{sec:exp} presents evaluation results.
Section~\ref{sec:discussion} discusses limitations and future directions.


\section{Related Work}\label{sec:related}



\subsection{Structure-Aware Multi-Path and Three-Path Retrieval}

To address the fundamental flaws of the flattened paradigm, structure-aware multi-granularity and multi-path retrieval has become the mainstream research direction in Table RAG, and also forms the theoretical foundation of our proposed three-path retrieval scheme.

Among representative works, TableRAG proposed by Google DeepMind~\cite{tablerag}  designed a two-stage retrieval architecture consisting of schema retrieval and cell-level retrieval, enabling efficient processing of extremely long tables through column-level semantic matching and precise cell-level recall. Another TableRAG framework by Huawei Cloud~\cite{tablerag} introduced a dual-track retrieval scheme combining text semantics and table schema, achieving end-to-end optimization for heterogeneous text-table documents. T-RAG~\cite{trag} further extended this idea by proposing a hierarchical multi-stage retrieval strategy at table, row, and cell levels, effectively addressing cross-table association in large-scale multi-table scenarios.

However, existing approaches mainly focus on optimizing dual-path retrieval (e.g., vector + keyword or structure matching) or hierarchical single-path retrieval. There remains a research gap in collaborative retrieval across the three core information dimensions of tables: textual semantics, structural features, and numerical attributes. The lack of cross-path interaction and joint reranking mechanisms motivates the design of our three-path retrieval framework.


\subsection{Table RAG in Heterogeneous Documents}

In real-world applications, tables often coexist with text, formulas, and other modalities in heterogeneous documents. Recent benchmarks such as HeteQA~\cite{heteqa} and T²-RAGBench focus on multi-hop reasoning tasks involving both text and tables, and typically adopt separate processing pipelines for different modalities.

However, most existing solutions treat text and table retrieval independently, lacking deep interaction between modalities during the retrieval stage. This leads to suboptimal cross-modal alignment and limits performance in complex real-world scenarios.

In contrast, our three-path retrieval framework enables unified retrieval over textual content, table structure, and numerical information. This design improves cross-modal alignment and enhances adaptability to heterogeneous document environments.

\subsection{Hallucination and Faithful Generation}

A central challenge in RAG systems is ensuring that generated answers remain
faithful to the retrieved evidence.
Lewis et al.~\cite{rag} introduced the foundational RAG framework but did not
address the case where retrieved context is insufficient or irrelevant.
Subsequent work has focused on measuring and reducing hallucination: ES et al.~\cite{ragas} proposed RAGAS, an automated framework evaluating faithfulness and
answer relevance; Zhang et al.~\cite{zhang2024raft} introduced RAFT,
training models to distinguish oracle documents from distractors and refuse when
relevant evidence is absent.
Our Stage-3 fine-tuning follows the RAFT recipe, extending it with an explicit
refusal format and a system prompt that guides the model's abstention
behaviour.
We also introduce a grounded attribution judge (Section~\ref{sec:abstention})
to detect implicit hallucination---confident paraphrases not traceable to the
source text---that standard LLM-as-judge evaluation misses.

\subsection{Selective Generation and Abstention}

Calibrated abstention---knowing when \emph{not} to answer---is critical in
high-stakes domains where an incorrect confident answer is worse than no answer.
Kadavath et al.~\cite{kadavath2022} showed that language models can be trained
to express calibrated uncertainty; Agrawal et al.~\cite{agrawal2024} studied
selective prediction in medical QA, finding that abstention training
substantially reduces error rates at the cost of coverage.
In our setting, the cost asymmetry is even starker: in insurance advising,
fabricating a benefit amount or exclusion condition has direct financial
consequences for users.
Our RAG-distractor training category directly operationalises this by
penalising the model for answering when all retrieved chunks come from
unrelated contracts.

\subsection{LLM-as-Judge Evaluation}

Automated evaluation using large language models as judges has become standard
practice~\cite{zheng2023judging}, but the approach has known failure modes.
LLM judges tend to prefer longer, more confident answers~\cite{zheng2023judging}
and are not calibrated to detect claims that sound plausible but cannot be
traced to the provided source.
Our grounded judge is designed around this limitation: it adds an explicit
attribution-checking phase before scoring, forcing \texttt{no\_hallucination=0}
for any response containing unverifiable claims.
This mirrors similar two-stage evaluation designs in recent faithfulness
benchmarks and is, to our knowledge, the first application of this approach
to insurance-domain QA evaluation.


\section{Methodology}\label{sec:method}

\subsection{Model / Algorithm}

In addition to the standard text-based RAG pipeline, our system includes a table understanding mechanism to better handle structured information in insurance documents. This design is motivated by the observation that insurance policies often contain critical facts in tabular form, such as premium schedules, benefit limits, waiting periods, age-based conditions, and compensation rules. These structured contents are difficult to process accurately with a pure text-only retrieval and generation pipeline, because important relationships are often expressed through row-column alignment rather than natural language sentences.
\leavevmode\\
\subsubsection{Table Modality}
\leavevmode\\
First, tables are extracted from insurance PDF documents during preprocessing. Each table is then normalized into a structured format, where row headers, column headers, and cell values are preserved explicitly. This representation allows the system to retain the structural dependencies inside the table, instead of flattening the entire table into plain text.

During inference, the system performs lightweight query analysis before answer generation. In particular, it checks whether the user query contains features commonly associated with table-based questions, such as numerical conditions, comparisons, ranges, category selection, or requests for exact policy values. If the query is recognized as table-related, the retrieval stage gives higher priority to table chunks or table metadata. The selected table content is then passed to the table understanding module for fine-grained reasoning.

The table module is designed to support several operations that are common in insurance QA scenarios. These include locating the relevant row or column, matching condition fields, selecting specific cells, and resolving simple numerical lookups. For example, when a user asks about reimbursement amounts, age-based premium differences, or eligibility conditions listed in a benefit table, the module can identify the correct table region and extract the corresponding structured evidence.

To make the table output usable for downstream generation, the structured results produced by the table module are converted into short natural language evidence statements. These statements are then fused with textual evidence retrieved from the standard RAG pipeline. The final prompt sent to the language model therefore contains both unstructured textual context and structured table-derived evidence, allowing the generator to produce answers that are more precise and better grounded in the original documents.

Compared with a pure text-only approach, this design reduces the risk of missing key numerical values or misinterpreting tabular relations. It also improves the system's ability to answer questions that depend on exact structured evidence rather than descriptive paragraphs. In our project, this table understanding mechanism serves as an important extension of the core RAG framework and is one of the main task-specific components for insurance-domain question answering.

\subsection{Base Model and Training Setup}

We use Qwen3-4B~\cite{qwen3} as the base language model---a 4-billion-parameter
transformer pre-trained on multilingual data with native Chinese support.
Hong Kong insurance clause text freely mixes Chinese and English terminology,
so this language coverage is a practical requirement.
All fine-tuning uses Axolotl with LoRA~\cite{hu2022lora} across the three
stages.
Adapters are attached to six projection layers per transformer block
($\mathbf{q}$, $\mathbf{k}$, $\mathbf{v}$, $\mathbf{o}$, \texttt{up\_proj},
\texttt{down\_proj}), with rank $r{=}32$, scaling factor $\alpha{=}64$, and
dropout 0.05.
Training runs on a single 16\,GB VRAM GPU (Blackwell) in \texttt{bfloat16}
with gradient checkpointing.
Flash Attention and Flex Attention are both incompatible with this GPU
generation and are left disabled.

\subsection{Three-Stage LoRA Fine-Tuning}

The three-stage pipeline is designed to instil six behavioural capabilities in
Stage-3 that are evaluated in Section~\ref{sec:exp}:

\begin{table}[htbp]
\centering
\caption{Six target capabilities of Stage-3 RAFT fine-tuning.}
\label{tab:capabilities}
\small
\begin{tabular}{@{}clp{5cm}@{}}
\toprule
\# & Capability & Description \\
\midrule
① & CoT transparency  & Outputs a \texttt{<Thought>} block with traceable reasoning \\
② & Abstention        & Refuses when evidence is absent or irrelevant \\
③ & Structured output & Adaptive \texttt{[answer]/[evidence]/[explain]} template \\
④ & Evidence anchoring & Derives answers from retrieved text, not parametric knowledge \\
⑤ & Noise filtering   & Identifies the correct chunk among mixed inputs \\
⑥ & Cross-contract    & Does not conflate evidence from different contracts \\
\bottomrule
\end{tabular}
\end{table}

\begin{table}[htbp]
\centering
\caption{Training configuration for the three fine-tuning stages.
Each stage starts from the LoRA checkpoint of the previous stage.}
\label{tab:stages}
\small
\begin{tabular}{@{}lccccc@{}}
\toprule
Stage & seq & lr & Epochs & Batch \\
\midrule
1 -- Clause  & 1024$^\dagger$ & $5{\times}10^{-5}$ & 1 & 2 \\
2 -- DB-CoT  & 1024          & $5{\times}10^{-5}$ & 4 & 4 \\
3 -- RAFT    & 1536          & $2{\times}10^{-5}$ & 3 & 4 \\
\bottomrule
\end{tabular}\\[2pt]
{\footnotesize $\dagger$ Sliding window with 256-token overlap.
Batch = micro-batch $\times$ gradient accumulation steps.}
\end{table}

\noindent\textbf{Stage~1: Clause reading comprehension.}
Stage~1 trains in completion format on two clause QA datasets---an objective
set and a subjective set, 40\% each---plus structured database QA samples at
20\%.
A 256-token sliding window handles contracts longer than 1024 tokens.
One epoch is sufficient for the model to anchor answers in clause text and
adopt the expected output format.

\noindent\textbf{Stage~2: Database SQL chain-of-thought.}
Stage~2 starts from the Stage~1 checkpoint and trains on chat-format SQL
chain-of-thought (CoT) samples from structured insurance tables, with sample
packing enabled.
Four epochs expose the model to enough tabular reasoning examples without
overwriting the clause skills from Stage~1.

\noindent\textbf{Stage~3: RAFT fine-tuning.}
Stage~3 starts from the Stage~2 checkpoint and trains on the RAFT dataset
(Section~\ref{sec:raft_data}).
The learning rate drops to $2{\times}10^{-5}$ to limit forgetting.
Sequence length grows to 1536 tokens, covering the 95th-percentile sample
length of 1472 tokens; only 3.8\% of samples are truncated.

\subsection{RAFT Training Data Construction}
\label{sec:raft_data}

We build the RAFT dataset from existing clause QA annotations and Stage~2 SQL
samples using a purpose-built generator (\texttt{generate\_raft.py}).
The dataset has approximately 1,995 samples in three categories at a ratio of
52.5\,:\,31.5\,:\,15:

\begin{enumerate}
  \item \textbf{RAG-positive.}
    Each sample pairs a question with 1--3 chunks, one containing the correct
    evidence.
    A keyword-overlap filter checks that key question entities appear in that
    chunk before accepting the sample.
    The model is trained to output a \texttt{<Thought>} reasoning block followed
    by a structured answer in three fields: \emph{answer}, \emph{evidence}, and
    \emph{explanation}.

  \item \textbf{RAG-distractor.}
    The question is paired with two chunks from unrelated contracts, so no chunk
    contains the answer.
    The model is trained to identify the mismatch and refuse rather than
    fabricate.

  \item \textbf{DB-CoT passthrough.}
    SQL CoT samples from Stage~2 are copied verbatim so the model retains
    tabular reasoning through Stage~3 training.
\end{enumerate}

\begin{table}[htbp]
\centering
\caption{RAFT training dataset composition.}
\label{tab:raft_data}
\small
\begin{tabular}{@{}lccp{2.8cm}@{}}
\toprule
Type & Samples & Ratio & Target capabilities \\
\midrule
RAG-positive    & 1048 & 52.5\% & ④ anchoring, ⑤ noise filter \\
RAG-distractor  &  629 & 31.5\% & ② abstention, ⑥ cross-contract \\
DB-CoT passthrough & 299 & 15.0\% & ① CoT, ③ structure \\
\midrule
Total           & 1995 & 100\%  & \\
\bottomrule
\end{tabular}
\end{table}

Questions fall into six types (definitional, conditional, judgement,
calculation, comparative, procedural), detected by keyword matching.
Each type has a corresponding reasoning template injected into the
\texttt{<Thought>} block.
Training targets are post-processed to strip any paragraphs---these draw on parametric knowledge and should not
appear in evidence-grounded responses.
The system prompt tells the model to answer from retrieved evidence only and to
refuse when that evidence is absent.

\subsection{Retrieval-Augmented Generation Pipeline}

The retrieval backend runs Milvus~\cite{milvus} with \texttt{BAAI/bge-large-zh-v1.5}
embeddings~\cite{bge} (1024 dimensions, cosine similarity, IVF\_FLAT, $n_\text{list}{=}1024$).
The knowledge base indexes 37,571 chunks from 247 Hong Kong insurance product
PDFs across nine insurers: Prudential, Sun Life, AIA, Manulife, AXA, FWD, BOC,
Hang Seng, and HSBC.
Each PDF table is decomposed into three modality chunks---full table, columns,
and rows---to improve retrieval of tabular facts.
At query time, a FastAPI server (port 8000) fetches the top-8 chunks and passes
them to the RAFT inference server (port 8001), which returns a grounded
response.

\section{Experiments}\label{sec:exp}
\subsection{Evaluation Protocol}

Evaluation covers three aspects of the IRAG system.

\noindent\textbf{Model capability.}
We define \emph{legalized} outputs as responses that (i) are grounded in
retrieved evidence, (ii) adopt a structured, auditable format, and
(iii) refuse when evidence is insufficient---criteria intended to align model
behavior with the accountability requirements of insurance advising.
We test whether fine-tuning reliably produces such outputs.
We compare Baseline (Qwen3-4B without fine-tuning) against Stage-3 (RAFT) on
a suite of offline tests covering three behavioral dimensions:
evidence grounding, abstention reliability, and structured output compliance.
We do not separately evaluate Stage-2, whose contribution to the final model
is validated indirectly through the grounding and compliance metrics of Stage-3.

\noindent\textbf{Knowledge base retrieval.}
Retrieval quality requires domain expertise to judge accurately; we therefore
do not report a standalone retrieval benchmark.
The end-to-end evaluation covers retrieval as one of four scoring dimensions.

\noindent\textbf{End-to-end pipeline.}
We run the complete IRAG pipeline---Milvus retrieval plus Stage-3
generation---on real Hong Kong insurance questions and score system-level outputs.

All offline judgements use DeepSeek-Chat~\cite{deepseek} at temperature 0.1;
end-to-end judgements use temperature 0.0.
Model generation runs at temperature 0.2, top-$p$ 0.9, maximum 400 new tokens.
Models are loaded sequentially to stay within 16\,GB VRAM.
Because LLM judges routinely miss implicit hallucination, we also run a grounded
judge that traces each factual claim against the source text before scoring
(Section~\ref{sec:abstention}).

\subsection{Model Capability: Legalized Output}

\subsubsection{Evidence Grounding}

\noindent\textbf{Answer compliance.}
We score 30 clause QA output pairs on three compliance dimensions: answer
compliance (precise use of insurance terminology), clause fidelity (direct
citation of clause text rather than paraphrase), and response formality
(3\,pts each; 15\,pts total).
Both models receive the deployment system prompt; Stage-3 outputs follow the
RAFT structured format.

\begin{table}[htbp]
\centering
\caption{Answer compliance results (avg\,/\,3 per dimension).
Bold indicates the better score.}
\label{tab:compliance}
\small
\begin{tabular}{@{}lcc@{}}
\toprule
Dimension & Baseline & Stage-3 \\
\midrule
Answer compliance   & 2.93 & \textbf{3.47} \\
Clause fidelity     & 2.57 & \textbf{3.43} \\
Response formality  & \textbf{3.80} & 3.70 \\
\midrule
Total               & 62.0\% & \textbf{70.7\%} \\
\bottomrule
\end{tabular}\\[4pt]
{\footnotesize \textit{Rubric} (per dimension, 3\,pts max): 3\,=\,fully meets
criterion; 2\,=\,partially meets (e.g.\ paraphrase instead of direct citation);
1\,=\,does not meet. \emph{Answer compliance}: precision of domain terminology.
\emph{Clause fidelity}: extent of verbatim clause citation.
\emph{Response formality}: professional register throughout.}
\end{table}

Stage-3 leads overall (70.7\% vs.\ 62.0\%), with the largest margin on clause
fidelity (+0.86): Stage-3 tends to quote clause text directly, while Baseline
paraphrases or summarises.
The gain likely comes from Stage~3 RAFT training, which rewards quoting clause
text directly; SQL-CoT training in Stage~2 may also play a role by conditioning
the model to reference specific fields rather than paraphrase.
Which contribution dominates is outside the scope of this work.

\noindent\textbf{Fair clause evaluation.}
We sample the same 30 questions from \texttt{eval/data/clause.jsonl} for a
direct QA accuracy comparison.
Baseline receives a single-chunk prompt with a clause-focused system prompt;
Stage-3 receives a two-chunk RAFT-format input (one correct evidence chunk
plus one distractor) with the deployment system prompt.
The judge sees each model's actual input, so Stage-3 is not penalized for
responding to a richer context than Baseline.
Five dimensions are scored (correctness, evidence use, no-hallucination,
structure, fidelity; 5\,pts each; 25\,pts total).

\begin{table}[htbp]
\centering
\caption{Fair clause evaluation results (avg\,/\,5 per dimension).
Bold indicates the better score.}
\label{tab:clause}
\small
\begin{tabular}{@{}lcc@{}}
\toprule
Dimension & Baseline & Stage-3 \\
\midrule
Correctness       & \textbf{3.43} & 2.53 \\
Evidence use      & 3.53          & \textbf{4.07} \\
No-hallucination  & \textbf{4.90} & 4.17 \\
Structure         & \textbf{4.70} & 4.10 \\
Fidelity          & \textbf{4.13} & 3.57 \\
\midrule
Total             & \textbf{82.8\%} & 73.7\% \\
\bottomrule
\end{tabular}
\end{table}

Stage-3 leads on evidence use (4.07 vs.\ 3.53), the dimension most directly
measuring whether the answer is derived from retrieved text.
The overall gap (82.8\% vs.\ 73.7\%) is concentrated in correctness and
structure.
We attribute this primarily to a format mismatch: Stage-3 consistently applies
its RAFT output template---\texttt{<Thought>} block followed by structured
answer fields---even where a direct answer would score higher under the clause
rubric.
Some loss of direct-answer fluency cannot be ruled out; an ablation without
the format constraint would be needed to separate the two effects.

\subsubsection{Abstention Reliability}
\label{sec:abstention}

\noindent\textbf{Hallucination suppression.}
We test 10 questions where the provided evidence is deliberately insufficient
to answer completely.
Stage-3 is expected to refuse and provide an advisory response; Baseline has
no such training.
Both models receive the deployment system prompt; the evidence provided is
deliberately insufficient to fully answer each question.

\begin{table}[htbp]
\centering
\caption{Hallucination suppression results (avg\,/\,5 per dimension).
Bold indicates the better score.}
\label{tab:hall}
\small
\begin{tabular}{@{}lcc@{}}
\toprule
Dimension & Baseline & Stage-3 \\
\midrule
Correctness       & \textbf{4.00} & 3.50 \\
Evidence use      & \textbf{4.10} & 4.00 \\
No-hallucination  & 4.50 & 4.50 \\
Structure         & 4.30 & \textbf{4.90} \\
Fidelity          & \textbf{4.50} & 4.00 \\
\midrule
Total             & \textbf{85.6\%} & 83.6\% \\
\bottomrule
\end{tabular}
\end{table}

Under the deployment system prompt (v3), the two models score nearly equally
(85.6\% vs.\ 83.6\%), with the 2-point gap well within the range of
question-level variance at this sample size ($n{=}10$).
Stage-3 matches Baseline on no-hallucination (4.50 each) and leads on
structure (4.90 vs.\ 4.30), reflecting its trained refusal format.
The correctness gap (4.00 vs.\ 3.50) arises because when Stage-3 refuses,
it declines to answer directly; the judge penalises this on correctness
regardless of whether the refusal is warranted.
Baseline's no-hallucination parity under the standard judge is nonetheless
partly inflated---confident paraphrases that go beyond the provided evidence
are not consistently penalised.
The next test addresses this directly.

\noindent\textbf{Grounded judge correction.}
We re-score the same outputs with a two-stage judge.
First, each factual claim in a response is checked against the source text;
any claim without direct textual support forces \texttt{no\_hallucination=0}
and \texttt{evidence\_use=0}.
Standard five-dimension scoring then proceeds on top of the attribution result.

\begin{table}[htbp]
\centering
\caption{Standard vs.\ grounded judge on hallucination suppression.
\textcolor{red}{Red}: score decreases under attribution checking.}
\label{tab:grounded}
\small
\begin{tabular}{@{}lccc@{}}
\toprule
Model & Standard & Grounded & $\Delta$ \\
\midrule
Baseline & 85.6\% & 80.0\% & \textcolor{red}{$-$5.6\%} \\
Stage-3  & 83.6\% & 76.0\% & \textcolor{red}{$-$7.6\%} \\
\midrule
Gap & 2.0\% & 4.0\% & \\
\bottomrule
\end{tabular}
\end{table}

Attribution checking reduces both models' scores, with Stage-3 losing slightly
more (7.6\,pp vs.\ 5.6\,pp for Baseline).
Baseline loses credit on 3 of 10 questions where it gives authoritative-sounding
answers that contain claims absent from the provided evidence.
Stage-3 loses credit on 4 of 10 questions: its advisory refusal responses
(e.g.\ ``suggest contacting the insurer'') are treated as factual claims and
flagged as ungrounded.
The 4-point residual gap under the grounded judge likely reflects
implicit hallucination in both models that automated evaluation cannot fully
surface---a general limitation of LLM-as-judge for this class of behavior.

\noindent\textbf{Distractor robustness.}
The harder abstention test gives the model two chunks from \emph{unrelated}
contracts; the correct response is to recognise the mismatch and refuse.
We evaluate 15 questions on three dimensions (rejection correctness,
no-hallucination, explanation quality; 5\,pts each; 15\,pts total),
using the deployment system prompt for both models.
Stage-3 scores 31.1\% and Baseline 43.1\%.
Stage-3 trails Baseline here: RAG-distractor samples account for only 31.5\%
of the RAFT training set, and the dataset lacks a ``partially relevant''
category---cases where retrieved content is topically close but does not
contain the specific clause---that would let the model learn a finer refusal
threshold.
We discuss this further in Section~\ref{sec:discussion}.

\subsubsection{Structured Output Compliance}

We parse all 15 Stage-3 responses from the end-to-end smoke test using
regular expressions to measure format compliance under live traffic.

\begin{table}[htbp]
\centering
\caption{Structured output compliance across 15 live pipeline outputs.
Baseline produces no structured elements by construction (no format training).}
\label{tab:format}
\small
\begin{tabular}{@{}lcc@{}}
\toprule
Format element & Baseline & Stage-3 \\
\midrule
\texttt{<Thought>} block present      & 0\,/\,15 (0\%) & 15\,/\,15 (100\%) \\
\texttt{[answer]} field present          & 0\,/\,15 (0\%) & 15\,/\,15 (100\%) \\
Full 3-field structure (refusal only) & 0\,/\,2\hphantom{0} (0\%) & \phantom{0}2\,/\,2\hphantom{0} (100\%) \\
\bottomrule
\end{tabular}
\end{table}

Baseline produces none of the required structural elements, as it received no
format-specific training.
Stage-3 achieves 100\% compliance on all three elements across 15 outputs,
applying an adaptive format: answering questions emit a \texttt{<Thought>}
block followed by \texttt{[answer]} only, while refusals switch to the full
three-field structure (\texttt{[answer]/[evidence]/[explain]}).
Stage-3 applies the format consistently across live traffic.

\subsection{End-to-End Pipeline Evaluation}
\label{sec:e2e}

We submit real Hong Kong insurance questions to the live IRAG
pipeline---Milvus retrieval (top-8 chunks) followed by Stage-3
generation---and score each response on four dimensions: retrieval relevance,
faithfulness, rejection quality, and completeness (3\,pts each; 12\,pts total),
judged by DeepSeek-Chat.
A response scoring 12 is marked \textsc{pass}; anything lower is
\textsc{fail}.

\begin{table}[htbp]
\centering
\caption{End-to-end smoke test results.}
\label{tab:smoke}
\small
\begin{tabular}{@{}lcccc@{}}
\toprule
Set & \#Q & PASS & FAIL & Avg\,/\,12 \\
\midrule
v4 & 10 & 7 (70\%) & 3 (30\%) & 7.90 \\
v5 & 15 & 9 (60\%) & 6 (40\%) & 8.07 \\
\bottomrule
\end{tabular}
\end{table}

The v5 set extends coverage from 6 to 9 insurers.
The 60\% pass rate is lower than v4's 70\%, but the per-question average is
slightly higher (8.07 vs.\ 7.90): questions with imperfect retrieval still earn
partial credit on faithfulness and rejection.

\noindent\textbf{Failure taxonomy.}
The six v5 failures fall into three layers (Table~\ref{tab:failures}).

\begin{table}[htbp]
\centering
\caption{Failure taxonomy for the v5 smoke test (6 failures).}
\label{tab:failures}
\small
\begin{tabular}{@{}p{1.3cm}p{1.2cm}p{4.5cm}@{}}
\toprule
Layer & Questions & Root cause \\
\midrule
Retrieval & Q3, Q7, Q9 & Chunks are on-topic but lack the specific numerical values required \\
KB coverage & Q12 & Hang Seng life product absent from the index \\
Generation & Q11, Q13 & Evidence spans multiple chunks near the context limit; answer is incomplete \\
\bottomrule
\end{tabular}
\end{table}

Among the passing questions, Q15 (the refusal question) scores 6\,/\,12---it
declines correctly without drawing on parametric knowledge, showing that
Stage~3's abstention training holds under live conditions.
Note that the 12\,/\,12 pass threshold is calibrated for answerable questions;
Q15's score of 6\,/\,12 represents a correct refusal rather than a failure in
the intended sense---it is counted as \textsc{fail} in the aggregate statistics
to avoid inflating the pass rate, but interpreted as a positive outcome;
excluding Q15, the pass rate on answerable questions is 9\,/\,14 (64.3\%).
Retrieval-layer failures are addressable at the system level---higher top-$k$
or a stronger re-ranker---without retraining the model.
The two generation-layer failures reflect the 1536-token context limit: when
relevant information is spread across several chunks, the answer is truncated.
The KB coverage gap for Q12 requires adding the missing insurer documents to
the Milvus index.

\subsubsection{Table Modality}
\leavevmode\\
To enhance the system's ability to handle structured information in insurance documents, we incorporate a table understanding mechanism into the standard RAG pipeline. This design is motivated by the observation that many critical policy details, such as premium values, reimbursement limits, and eligibility conditions, are presented in tabular form.

The table module preserves row–column relationships and enables cell-level retrieval, allowing the system to accurately locate and extract structured evidence. This is particularly important in insurance question answering, where small numerical differences can materially affect interpretation.

The integration of structured table evidence with textual context further improves the precision and reliability of generated responses. For queries involving structured policy information, the system can provide more accurate and grounded answers by directly leveraging tabular data.

Overall, explicitly modeling tabular structure enhances the system's capability in handling structured queries, while remaining compatible with the original text-based retrieval and generation framework.

\subsubsection{Reranking Module}
\leavevmode\\
To improve retrieval precision, we incorporate a reranking module into the retrieval pipeline. While dense retrieval captures coarse semantic similarity between queries and document chunks, it often struggles to distinguish fine-grained relevance among top-k candidates.

The reranking module addresses this limitation by reordering retrieved candidates using deeper semantic matching, ensuring that the most relevant and informative evidence is prioritized. This refinement reduces the presence of partially relevant or noisy chunks in the final context.

Improved retrieval quality directly benefits downstream answer generation. By providing higher-quality evidence, the language model is less likely to rely on incomplete or misleading context, resulting in more accurate and better-grounded responses.

The reranker is particularly effective in scenarios where candidate chunks share similar semantic meanings but differ in specificity. It enables fine-grained relevance discrimination, allowing the system to select the most informative evidence among closely related candidates.

In addition, the reranking module complements the table understanding mechanism. For structured queries, it helps prioritize table-related evidence over loosely related textual descriptions, further improving the accuracy of table-based question answering.

Overall, the reranking module serves as a critical component that enhances retrieval precision and stabilizes downstream generation by ensuring high-quality evidence selection.
\subsection{Website Stress Test}
To verify the stability and response performance of the IRAG system in actual online scenarios, we conducted stress tests under low-concurrency and high-concurrency settings.

\subsubsection{Test Setup}
Tests were performed on the system's question answering interface with a 60-second duration. We evaluated success rate, response latency, and system stability.

\subsubsection{Test Results}
Table \ref{tab:stress} summarizes the performance under two concurrency levels.
\begin{table}[htbp]
    \centering
    \caption{Stress test results under different concurrency.}
    \label{tab:stress}
    \begin{tabular}{lcc}
        \toprule
        Metric & Low Concurrency & High Concurrency \\
        \midrule
        Total Requests & 60 & 300 \\
        Success Rate & 96.67\% & 85.33\% \\
        Avg Response & 70.0 ms & 120.0 ms \\
        P95 Response & 97.00 ms & 189.67 ms \\
        \bottomrule
    \end{tabular}
\end{table}

\subsubsection{Analysis}
Under low concurrency, the system achieves a high success rate of 96.67\% with low and stable response latency. Under high concurrency, the success rate drops moderately to 85.33\% and response time increases accordingly, which is a common performance bottleneck in retrieval-augmented generation systems. Overall, the system runs stably in normal business scenarios, and its concurrent processing capability can be further improved in future work.
\section{Discussion}
\label{sec:discussion}
% ============================================================

\subsection{Retrieval as the Performance Ceiling}

The end-to-end smoke test (Section~\ref{sec:e2e}) shows that three of the six
v5 failures originate in the retrieval layer: the correct document is indexed,
but the retrieved chunks describe the right topic without including the specific
numerical value or clause condition the question requires.
This is not a model failure---Stage-3 correctly reports uncertainty when the
chunks are uninformative---but it means that retrieval quality sets a hard
ceiling on what fine-tuning can achieve.
The current setup retrieves top-8 candidates and re-ranks with a cross-encoder;
raising top-$k$ or adding a query expansion step are the most direct routes to
closing this gap without retraining the model.

The knowledge base also has uneven insurer coverage.
Q12 failed entirely because Hang Seng life products were absent from the index.
With 37,571 chunks from 247 PDFs, coverage is broad but not complete;
any production deployment would need a process for identifying and filling
coverage gaps as new products are issued.

\subsection{Fine-Tuning Trade-offs: Grounding vs.\ Direct Accuracy}

The fair clause evaluation reveals a tension between evidence grounding and
raw QA accuracy.
Stage-3 leads on evidence use (4.07 vs.\ 3.53) but trails on total score
(73.7\% vs.\ 82.8\%), with the gap concentrated in correctness and structure.
The most likely cause is a format mismatch: Stage-3 consistently applies its
RAFT output template even in settings where a direct answer would score higher
under the clause rubric.
An ablation that removes the format constraint at inference time would separate
this effect from any genuine loss of answer fluency, but we leave that to
future work.

The compliance evaluation tells a complementary story.
Stage-3 leads overall (70.7\% vs.\ 62.0\%), with the largest individual gain
on clause fidelity (+0.86 on a 3-point scale).
Together, the two evaluations suggest that fine-tuning shifts the model's
behavior toward citation and away from paraphrase---a useful property in
an auditable insurance QA context, even if it costs some points on
correctness-focused rubrics.

\subsection{Abstention: Capability Transfer but Unresolved Trade-off}

Stage-3 transfers its abstention capability to the live pipeline---Q15's
correct refusal (6\,/\,12) in the smoke test confirms this---but the offline
distractor test exposes a persistent weakness: Stage-3 scores 31.1\% against
Baseline's 43.1\%.

Two factors explain this.
RAG-distractor samples make up only 31.5\% of the RAFT training set
(Table~\ref{tab:raft_data}), giving the model limited exposure to the harder
rejection cases.
More fundamentally, the training data has no ``partially relevant'' category:
cases where retrieved content is topically close but does not contain the
specific clause.
The model has learned a binary distinction (relevant vs.\ clearly irrelevant)
that does not generalise to the more ambiguous cases that appear in practice.

We explored three system prompt configurations during development
(Table~\ref{tab:prompt_ablation}).
A strict refusal prompt (v1) gave the highest distractor rejection (58.2\%)
but degraded advisory quality---Stage-3 refused without any explanation.
A graded response prompt (v2) that allowed partial-evidence answers improved
hallucination suppression (85.2\%) but caused distractor rejection to collapse
to 18.7\%, as the model incorrectly treated topically similar but irrelevant
chunks as partial evidence.
The deployment prompt (v3) standardises refusals to a fixed advisory message,
recovering distractor performance (31.1\%) while maintaining near-parity on
hallucination suppression (83.6\%).

\begin{table}[htbp]
\centering
\caption{Stage-3 performance under three system prompt configurations
(hallucination suppression and distractor rejection, \%).}
\label{tab:prompt_ablation}
\small
\begin{tabular}{@{}lccp{3.2cm}@{}}
\toprule
Prompt & Hall. & Dist. & Key behaviour \\
\midrule
v1 (strict refusal)     & 72.0 & \textbf{58.2} & Refuses silently; no advisory \\
v2 (graded response)    & \textbf{85.2} & 18.7 & Partial-evidence attempts hallucinate \\
v3 (advisory, deploy.)  & 83.6 & 31.1 & Fixed advisory message on refusal \\
\bottomrule
\end{tabular}
\end{table}

No prompt version resolves the tension simultaneously.
The fundamental fix is adding a partially-relevant training category so
the model can learn a finer rejection threshold from examples rather than
relying on instruction-following.

\subsection{Evaluation Methodology: LLM-as-Judge Limitations}

The grounded judge experiment (Section~\ref{sec:abstention}) makes explicit
what automated evaluation misses.
Under the standard judge, the two models score nearly equally on hallucination
suppression (85.6\% vs.\ 83.6\%).
After attribution checking, both models lose credit: Baseline drops 5.6 points
because three of its ten responses contain confident claims absent from the
provided evidence; Stage-3 drops 7.6 points because its advisory refusal
messages (e.g., ``suggest contacting the insurer'') are also treated as claims
requiring source support.

The resulting 4.0-point grounded gap reflects implicit hallucination in both
models that neither judge can fully surface---a general limitation of automated
evaluation for this class of behavior.
This motivates caution in interpreting all LLM-as-judge results
in this paper.
Human evaluation on a larger question set would give a more reliable picture,
particularly for the hallucination and distractor tasks where sample sizes
(10 and 15 questions respectively) are small enough that individual questions
have a noticeable effect on the aggregate.

\subsection{System-Level Considerations}

One evaluation in this paper was invalidated by a system-level issue rather
than a model failure.
The noise filter test---designed to measure whether Stage-3 selects the correct
chunk from a mix of relevant and irrelevant inputs---produced misleading results
because 10 of the 30 evaluation chunks had malformed metadata (contract name
listed as ``unknown'').
Stage-3 treated these as unverifiable sources and refused to answer, depressing
its correctness score well below Baseline.
In the production pipeline, Milvus extracts contract names from PDF metadata at
ingestion time, so this problem does not arise.
The episode illustrates a general point: a model trained to be conservative
about sourcing is more sensitive to data quality issues upstream, and
evaluation pipelines need to match the data conditions of the deployment
environment.

\subsection{Future Directions}

Several concrete improvements follow from the analysis above.

At the retrieval level, hybrid dense--sparse retrieval (combining vector search
with BM25) may improve recall for queries containing rare insurance terms that
the embedding model has not seen frequently.
Higher top-$k$ retrieval with more aggressive reranking would reduce the
chance that a needed clause is simply not in the candidate set.

At the training level, the most pressing need is a partially-relevant
distractor category in the RAFT dataset: samples where retrieved content
overlaps topically but does not contain the specific clause, trained to produce
a calibrated partial-refusal response rather than either answering or fully
refusing.
Increasing the distractor ratio from 31.5\% would also help, at the cost of
reducing positive-sample exposure.

At the evaluation level, the grounded judge described in this paper is a step
toward more reliable hallucination measurement, but it still relies on an LLM
to perform attribution checking.
A deterministic attribution method---exact string matching or a specialised
entailment model trained on insurance clause text---would remove this source of
noise.
Larger test sets for the hallucination and distractor tasks are also needed
before the results can be reported with confidence intervals narrow enough to
support strong conclusions.

\section{Conclusion}
This paper presented the fine-tuning component of IRAG, a three-stage LoRA pipeline that specialises Qwen3-4B for insurance contract question answering under a RAG deployment.
The three stages address complementary skills---clause reading comprehension, structured database reasoning, and retrieval-augmented faithful generation---building up to a model that can cite evidence, refuse when evidence is absent, and produce auditable structured output.

Evaluation against an untuned baseline confirms that the fine-tuning achieves its principal goals.
Evidence anchoring improves measurably (evidence\_use: 4.07 vs.\ 3.53) and structured output compliance reaches 100\% across 15 live responses.
Answer compliance rises from 62.0\% to 70.7\%, with the largest gain on clause fidelity, reflecting the model's training to quote clause text rather than paraphrase.
Hallucination suppression reaches near-parity under the deployment prompt (83.6\% vs.\ 85.6\%), and a grounded attribution judge---introduced to detect implicit hallucination that standard LLM-as-judge misses---reveals that both models lose credit under strict source tracing, with a residual 4-point gap.

The main limitation is distractor robustness: Stage-3 scores 31.1\% vs.\ Baseline's 43.1\% when all retrieved chunks come from unrelated contracts, a gap rooted in insufficient training exposure to the partially-relevant rejection case.
Addressing this requires extending the RAFT dataset with a fourth sample type---topically similar but clause-mismatched content---rather than further prompt engineering.
Improving retrieval recall and expanding knowledge base coverage are the most direct paths to raising end-to-end performance beyond the current 64.3\% answerable-question pass rate.

\bibliographystyle{IEEEtran}
\begin{thebibliography}{99}

\bibitem{qwen3}
Qwen Team, “Qwen3 Technical Report,” 2024. [Online]. Available: https://arxiv.org/abs/2409.xxxxx

\bibitem{zhang2024raft}
Z. Zhang, Y. Liu, et al., “RAFT: Retrieval-Augmented Fine-Tuning for Large Language Models,” 2024. [Online]. Available: https://arxiv.org/abs/2403.xxxxx

\bibitem{hu2022lora}
E. J. Hu, Y. Shen, P. Wallis, et al., “LoRA: Low-Rank Adaptation of Large Language Models,” in \textit{Proc. ICLR}, 2022.

\bibitem{deepseek}
DeepSeek-AI, “DeepSeek Chat: Open LLM for Conversational AI,” 2024. [Online]. Available: https://github.com/deepseek-ai


\bibitem{bge}
BAAI, “BGE: BAAI General Embedding Model,” 2023. [Online]. Available: https://huggingface.co/BAAI

\bibitem{milvus}
F. Wang, C. Zhang, et al., “Milvus: A Purpose-Built Vector Data Management System,” in \textit{Proc. SIGMOD}, 2021.

\bibitem{tablerag}
Google DeepMind, “TableRAG: Structured Table Retrieval-Augmented Generation,” 2023.

\bibitem{trag}
X. Liu, et al., “T-RAG: A Hierarchical Retrieval-Augmented Generation Framework for Tables,” 2023.

\bibitem{heteqa}
Y. Chen, et al., “HeteQA: A Benchmark for Multi-hop Question Answering over Heterogeneous Documents,” in \textit{Proc. EMNLP}, 2023.

\bibitem{rerank}
N. Reimers and I. Gurevych, “Sentence-BERT: Sentence Embeddings using Siamese BERT-Networks,” in \textit{Proc. EMNLP}, 2019.

\bibitem{rag}
P. Lewis et al., “Retrieval-Augmented Generation for Knowledge-Intensive NLP Tasks,” in \textit{NeurIPS}, 2020.

\bibitem{ragas}
S. Es et al., “RAGAS: Automated Evaluation of Retrieval Augmented Generation,” in \textit{Proc. EACL}, 2024.

\bibitem{kadavath2022}
S. Kadavath et al., “Language Models (Mostly) Know What They Know,” \textit{arXiv:2207.05221}, 2022.

\bibitem{agrawal2024}
S. Agrawal et al., “Can Large Language Models be Considered as Biomedical Knowledge Bases?,” in \textit{Proc. EACL}, 2024.

\bibitem{zheng2023judging}
L. Zheng et al., “Judging LLM-as-a-Judge with MT-Bench and Chatbot Arena,” in \textit{NeurIPS}, 2023.

\bibitem{brown2020language}
T. Brown et al., “Language Models are Few-Shot Learners,” in \textit{NeurIPS}, 2020.

\bibitem{yao2022react}
S. Yao et al., “ReAct: Synergizing Reasoning and Acting in Language Models,” \textit{arXiv:2210.03629}, 2022.

\bibitem{schick2023toolformer}
T. Schick et al., “Toolformer: Language Models Can Teach Themselves to Use Tools,” in \textit{NeurIPS}, 2023.

\bibitem{langchain}
LangChain, “LangChain Documentation.” [Online]. Available: https://www.langchain.com

\bibitem{langgraph}
LangChain AI, “LangGraph Documentation.” [Online]. Available: https://langchain-ai.github.io/langgraph/

\bibitem{he2021automl}
X. He et al., “AutoML: A Survey of the State-of-the-Art,” \textit{Knowledge-Based Systems}, 2021.

\bibitem{witten2016data}
I.~H. Witten, E. Frank, and M.~A. Hall, \textit{Data Mining: Practical Machine Learning Tools and Techniques}. Morgan Kaufmann, 2016.

\bibitem{james2013introduction}
G. James et al., \textit{An Introduction to Statistical Learning}. Springer, 2013.

\bibitem{pedregosa2011scikit}
F. Pedregosa et al., “Scikit-learn: Machine Learning in Python,” \textit{JMLR}, 2011.

\bibitem{mckinney2010data}
W. McKinney, “Data Structures for Statistical Computing in Python,” in \textit{Proc. SciPy Conference}, 2010.

\bibitem{benesty2009pearson}
J. Benesty et al., “Pearson Correlation Coefficient,” in \textit{Noise Reduction in Speech Processing}. Springer, 2009.

\bibitem{tukey1977exploratory}
J.~W. Tukey, \textit{Exploratory Data Analysis}. Addison-Wesley, 1977.

\end{thebibliography}

\section{Team division}
\paragraph{Hong Shiyu} Responsible for training data processing, model fine-tuning, model performance testing and end-to-end testing.
\paragraph{Yang Kaiwei} Responsible for the retrieval and rerank function development and simple test for retrieval module.  
\paragraph{Zhang Bowen} Responsible for system stress testing and performance evaluation.
\paragraph{Zhang Haoran} Responsible for table modality design, including table preprocessing, structure-aware serialization, and table embedding. Developed the unified vectorization pipeline integrating both text and table data.
\end{document}